\chapter{Sitio público}
\label{cap:sitio-publico}

\begin{procedimiento}{Objetivo del capítulo}{Visitante y participante}
Explicar cómo consultar la información pública del torneo, revisar reglas y patrocinadores, e ingresar al proceso de registro.
\end{procedimiento}

\section{Página de inicio}

La página de inicio presenta el nombre del torneo, el estado de la edición activa cuando existe, conteos de colegios y universidades, reglas destacadas y accesos principales. Desde esta pantalla el usuario puede seleccionar \botoninterfaz{Registrar un robot} o \botoninterfaz{Consultar reglamento}. La figura de referencia para esta pantalla se presentó en la \cref{fig:bloque1-mu-001}.

\begin{procedimiento}{Consultar información pública del torneo}{Visitante}
Permite revisar la información general antes de iniciar una inscripción.
\end{procedimiento}

\paso{1}{Abra la página de inicio del torneo.}
\paso{2}{Revise el bloque \campointerfaz{Estado del evento} para identificar si hay edición activa.}
\paso{3}{Use \botoninterfaz{Consultar reglamento} si desea revisar reglas antes de registrarse.}
\paso{4}{Use \botoninterfaz{Registrar un robot} cuando quiera iniciar el registro.}

\begin{resultadoesperado}
El visitante identifica el evento, encuentra las rutas de consulta y sabe dónde comenzar el registro.
\end{resultadoesperado}

\section{Reglas del torneo}

La pantalla de reglas muestra una consulta rápida, la edición activa cuando está disponible, el botón \botoninterfaz{Abrir PDF completo} y las secciones de reglas publicadas. El PDF se abre en una pestaña o visor del navegador, según la configuración del equipo.

\figuraManual{capturas/finales/MU-002-reglas-publicas.png}{Página de reglas donde el visitante puede revisar recomendaciones previas y abrir el reglamento oficial en PDF.}{fig:bloque1-mu-002}

\begin{procedimiento}{Consultar reglas oficiales}{Visitante o participante}
Use este procedimiento antes de registrar un robot o antes de asistir al torneo.
\end{procedimiento}

\paso{1}{Seleccione \botoninterfaz{Reglas} en el encabezado público o \botoninterfaz{Consultar reglamento} desde la página de inicio.}
\paso{2}{Lea la sección \campointerfaz{Antes de registrarte}.}
\paso{3}{Presione \botoninterfaz{Abrir PDF completo} si necesita revisar el documento oficial completo.}
\paso{4}{Regrese al sitio público usando el navegador o los enlaces del encabezado.}

\begin{fallas}
Si el PDF no carga, revise la conexión y vuelva a intentarlo. Si el problema continúa, consulte con la organización antes de enviar un registro, porque el reglamento puede contener condiciones obligatorias para participar.
\end{fallas}

\section{Patrocinadores}

La pantalla de patrocinadores muestra los aliados del torneo. Algunas tarjetas pueden abrir el sitio oficial del patrocinador; otras pueden mostrar el mensaje \campointerfaz{Enlace oficial en confirmación} cuando el enlace aún no está disponible.

\figuraManual{capturas/finales/MU-003-patrocinadores.png}{Galería pública de patrocinadores y aliados del torneo, con tarjetas informativas y enlaces cuando están disponibles.}{fig:bloque1-mu-003}

\begin{procedimiento}{Consultar patrocinadores}{Visitante}
Permite reconocer los aliados visibles en el sitio público.
\end{procedimiento}

\paso{1}{Seleccione \botoninterfaz{Patrocinadores} en el encabezado público.}
\paso{2}{Revise las tarjetas disponibles.}
\paso{3}{Si una tarjeta indica \campointerfaz{Visitar patrocinador}, puede abrir el enlace externo.}
\paso{4}{Si la tarjeta indica que el enlace está en confirmación, úsela solo como información visual.}

\begin{nota}
Los enlaces externos dependen de la información publicada por la organización. El manual no define ni valida contenido de sitios externos.
\end{nota}

\section{Acceso al registro desde el sitio público}

El sitio público ofrece varias entradas al registro: el enlace \botoninterfaz{Registro} del encabezado, el botón \botoninterfaz{Registrar un robot} de la página de inicio y los llamados a la acción asociados a la inscripción. El procedimiento completo de registro se explica en el \cref{cap:registro-participantes}.
